\chapter{Background}\label{ch:background}

\paragraph{SOC estimation.} Coulomb counting integrates the current; it is cheap but drifts with sensor offset and noise. Reading the SOC from the open-circuit voltage (OCV) works only after the cell has rested. Model-based estimators combine both: an extended Kalman filter (EKF) predicts the cell state with an equivalent circuit model and corrects it with the measured voltage \citep{plett2004ekf}. The model used here is a resistance $R_0$ in series with two resistor--capacitor branches; all resistances depend on temperature through the Arrhenius law \citep{plett2015}. A sensor offset can be added to the filter state, but it is observable only where the OCV curve is curved \citep{zhao2017}. Data-driven estimators such as recurrent neural networks can learn the non-linear behaviour from history \citep{nejakar2025}.

\paragraph{Thermal behaviour.} Heat generated in the cell is $q = i\,(\mathrm{OCV} - v)$ \citep{bernardi1985}. A cylindrical cell is modelled by two thermal nodes, core and surface \citep{lin2014}. Only the surface can be measured, so the core temperature has to be estimated.

\paragraph{Digital twins and cloud BMS.} A digital twin mirrors the battery's state in software. Li et al.\ ran SOC and state-of-health estimation in the cloud \citep{li2020}; Nejakar et al.\ synchronised an ESP32 monitor with a cloud database and a dashboard \citep{nejakar2025}. Connecting a BMS to the cloud also opens security risks \citep{naseri2023}, so safety decisions should stay on the vehicle.

\paragraph{Hardware-in-the-loop testing.} In hardware-in-the-loop (HIL) testing a real controller runs against a simulated plant. Laboratory BMS rigs use real-time simulators and cell-voltage emulators \citep{barreras2016,bui2019}. When only the control strategy is tested, it is sufficient to simulate the cell modules over CAN \citep{haupt2013}, which is the approach taken here.
